\begin{tikzpicture}[
    scale=4.2,
    every node/.style={font=\small},
    spring/.style={decorate, decoration={zigzag, segment length=2.2pt, amplitude=2pt, pre length=2pt, post length=2pt}},
    arrowstyle/.style={-Stealth, thick},
    flowarrow/.style={-Stealth, thick},
    quotestyle/.style={<->, thin, gray!70},
    dashedaxis/.style={dash dot, thin, gray!60},
    refchord/.style={dashed, thin, black!55},
    leader/.style={thin, gray!70}
]

% ──────────────────────────────────────────────────────────────────
% Pose
\def\alphaDeg{-7}
\def\deltaDeg{-18}

% Geometry references (chord = 1)
\def\EAx{0.30}
\def\hingeX{0.75}

% ──────────────────────────────────────────────────────────────────
% NACA 2312 — main body (0 → 75% chord)
\def\mainAirfoil{
    (0.0000,+0.0000) (0.0015,+0.0071) (0.0062,+0.0143) (0.0138,+0.0216)
    (0.0245,+0.0290) (0.0381,+0.0363) (0.0545,+0.0434) (0.0737,+0.0502)
    (0.0955,+0.0566) (0.1198,+0.0624) (0.1464,+0.0675) (0.1753,+0.0719)
    (0.2061,+0.0753) (0.2388,+0.0778) (0.2730,+0.0791) (0.3087,+0.0793)
    (0.3455,+0.0788) (0.3833,+0.0775) (0.4218,+0.0757) (0.4608,+0.0732)
    (0.5000,+0.0703) (0.5392,+0.0668) (0.5782,+0.0630) (0.6167,+0.0588)
    (0.6545,+0.0544) (0.6913,+0.0498) (0.7270,+0.0450) (0.7500,+0.0418)
    (0.7500,-0.0183) (0.7270,-0.0199) (0.6913,-0.0223) (0.6545,-0.0247)
    (0.6167,-0.0270) (0.5782,-0.0293) (0.5392,-0.0315) (0.5000,-0.0335)
    (0.4608,-0.0353) (0.4218,-0.0369) (0.3833,-0.0381) (0.3455,-0.0389)
    (0.3087,-0.0394) (0.2730,-0.0394) (0.2388,-0.0394) (0.2061,-0.0392)
    (0.1753,-0.0388) (0.1464,-0.0380) (0.1198,-0.0368) (0.0955,-0.0351)
    (0.0737,-0.0329) (0.0545,-0.0301) (0.0381,-0.0267) (0.0245,-0.0227)
    (0.0138,-0.0180) (0.0062,-0.0127) (0.0015,-0.0067) (0.0000,+0.0000)
}

% NACA 2312 — flap upper surface (75% → TE)
\def\flapUpper{
    (0.7500,+0.0418) (0.7612,+0.0401) (0.7939,+0.0353) (0.8247,+0.0305)
    (0.8536,+0.0258) (0.8802,+0.0213) (0.9045,+0.0171) (0.9263,+0.0132)
    (0.9455,+0.0097) (0.9619,+0.0066) (0.9755,+0.0040) (0.9862,+0.0019)
    (0.9938,+0.0004) (0.9985,-0.0005) (1.0000,-0.0008) (1.0000,+0.0008)
    (0.9985,+0.0007) (0.9938,+0.0003) (0.9862,-0.0003) (0.9755,-0.0012)
    (0.9619,-0.0023) (0.9455,-0.0037) (0.9263,-0.0052) (0.9045,-0.0069)
    (0.8802,-0.0088) (0.8536,-0.0108) (0.8247,-0.0130) (0.7939,-0.0152)
    (0.7612,-0.0175) (0.7500,-0.0183)
}

% ──────────────────────────────────────────────────────────────────
% Freestream reference axis (horizontal, alpha reference)
\draw[dashedaxis] (-0.95, 0) -- (1.55, 0);

% ──────────────────────────────────────────────────────────────────
% Ceiling
\draw[thick] (0.10, 0.65) -- (0.55, 0.65);
\foreach \x in {0.12, 0.18, 0.24, 0.30, 0.36, 0.42, 0.48, 0.54}{
    \draw[thin] (\x, 0.65) -- (\x - 0.04, 0.72);
}

% ──────────────────────────────────────────────────────────────────
% MAIN AIRFOIL — rotated by alpha around the EA
\begin{scope}[rotate around={\alphaDeg:(\EAx, 0)}]
    \draw[thick, fill=black!4] plot coordinates \mainAirfoil -- cycle;

    % Chord line of the main airfoil — extended well past LE on the left so
    % that it serves as the reference line for the angle of attack alpha,
    % and past the hinge on the right.
    \draw[refchord] (-0.55, 0) -- (1.10, 0);

    % Elastic axis
    \filldraw[black] (\EAx, 0) circle (0.012);

    % Wing CG
    \coordinate (CGw) at (0.3, 0);
    \draw[thick] (CGw) circle (0.024);
    \fill[black] (CGw) ++(0:0.024) arc (0:90:0.024) -- (CGw) -- cycle;
    \fill[black] (CGw) ++(180:0.024) arc (180:270:0.024) -- (CGw) -- cycle;

    % Hinge (on main airfoil side, at 75% chord)
    \filldraw[black] (\hingeX, 0) circle (0.012);

    % Heave spring upper attachment
    \coordinate (Sattach) at (\EAx, 0.075);

    % Reference points exported to outer frame
    \coordinate (EArot)    at (\EAx, 0);
    \coordinate (Hingerot) at (\hingeX, 0);
    \coordinate (LErot)    at (0, 0);
    \coordinate (TEextrot) at (1.10, 0);
\end{scope}

% ──────────────────────────────────────────────────────────────────
% FLAP — rotated by alpha (with the wing) and then by delta around the hinge.
% LE of the flap is rounded with a semicircle joining upper and lower surfaces.
\begin{scope}[rotate around={\alphaDeg:(\EAx, 0)}]
    \begin{scope}[rotate around={\deltaDeg:(\hingeX, 0)}]
        % Flap outline: upper surface (LE→TE→back to LE_lower) + rounded LE arc
        \draw[thick, fill=black!4]
            plot coordinates \flapUpper
            arc (270 : 90 : 0.0301)   % rounded LE: semicircle from (0.75,-0.0183) to (0.75,+0.0418)
            -- cycle;
        % NOTE on the arc:
        % center of the semicircle = midpoint of the two LE points = (0.75, +0.01175)
        % radius                   = (0.0418 - (-0.0183)) / 2 = 0.03005
        % drawn from angle 270° (bottom) to 90° (top), bulging to the LEFT (towards main wing)

        % Hinge point at the centre of the rounded LE of the flap
        \filldraw[black] (0.75, 0.01175) circle (0.012);

        % Flap chord line (hinge to TE), dashed reference for delta
        \draw[refchord] (\hingeX, 0) -- (1.10, 0);

        % Flap CG
        \coordinate (CGf) at (0.87, 0);
        \draw[thick] (CGf) circle (0.020);
        \fill[black] (CGf) ++(0:0.020) arc (0:90:0.020) -- (CGf) -- cycle;
        \fill[black] (CGf) ++(180:0.020) arc (180:270:0.020) -- (CGf) -- cycle;
    \end{scope}
\end{scope}

% ──────────────────────────────────────────────────────────────────
% HEAVE SPRING k_h
\draw[spring, thick] (\EAx, 0.65) -- (Sattach);
\node[right] at (\EAx + 0.025, 0.45) {$k_{h}$, $d_{h}$};

% ──────────────────────────────────────────────────────────────────
% TORSIONAL SPRING k_alpha — compact spiral on the EA
\begin{scope}[shift={(EArot)}]
    \draw[thick]
        (0:0.035) arc (0:330:0.035)
        (330:0.035) -- (330:0.050)
        arc (330:30:0.050)
        (30:0.050) -- (30:0.065)
        arc (30:330:0.065);
\end{scope}
% Label k_alpha placed UP-LEFT, near the heave spring, with a leader line
\node at (\EAx - 0.18, 0.32) {$k_{\alpha}$, $d_{\alpha}$};
\draw[leader] (\EAx - 0.16, 0.30) -- (\EAx - 0.04, 0.05);

% ──────────────────────────────────────────────────────────────────
% Labels EA, hinge, CG_w, CG_f
% EA label shifted to the right so it does not overlap with the heave spring
\node[above] at (\EAx + 0.11, 0.15) {EA};
\draw[leader] (\EAx + 0.10, 0.15) -- (\EAx + 0.015, 0.03);

\node[above] at (0.78, 0.40) {hinge};
\draw[leader] (0.78, 0.38) -- (\hingeX + 0.005, 0.05);

\node[below] at (0.42, -0.45) {CG$_{w}$};
\draw[leader] (0.42, -0.43) -- (CGw);

% CG_f label with leader pointing to the actual CG_f marker.
% CGf in the rotated (alpha + delta) frame is at (0.87, 0); in world frame
% this maps to approximately (0.855, -0.106).
\node[below] at (0.95, -0.32) {CG$_{f}$};
\draw[leader] (0.93, -0.30) -- (0.858, -0.115);

% ──────────────────────────────────────────────────────────────────
% FREESTREAM U_inf — placed UP, well separated from h
\draw[flowarrow] (-0.85, 0.25) -- (-0.55, 0.25);
\node[above] at (-0.70, 0.26) {$U_{\infty}$};

% ──────────────────────────────────────────────────────────────────
% PLUNGE h — measured from the horizontal freestream reference axis
% (y = 0), which is the same dash-dot line that also defines the vertical
% offset d_y on the right side of the figure. No extra dashed reference
% line is needed: h is simply a downward arrow starting on that axis.
\draw[arrowstyle] (-0.85, 0) -- (-0.85, -0.30);
\node[left] at (-0.87, -0.15) {$h(t)$};

% ──────────────────────────────────────────────────────────────────
% ANGLE OF ATTACK alpha — defined with the SAME LOGIC as delta:
% a compact arc drawn at the vertex (the LE) between the zero-lift
% line (airfoil chord) and the horizontal reference axis that passes
% through the EA. Drawn inside the alpha-rotated scope so the chord
% sits at angle 0° (local x-axis); the world-horizontal line then
% sits at angle -alphaDeg in this local frame.
\begin{scope}[rotate around={\alphaDeg:(\EAx, 0)}]
    \draw[thick] (-0.22, 0) arc (180 : 180 - \alphaDeg : 0.52);
    \node at (-0.30, 0.055) {$\alpha(t)$};
\end{scope}

% ──────────────────────────────────────────────────────────────────
% FLAP DEFLECTION delta — angle between the extension of the main chord
% and the flap chord. Drawn ABOVE the hinge for clarity, ample arc.
\begin{scope}[rotate around={\alphaDeg:(\EAx, 0)}]
    \draw[thick] (\hingeX + 0.22, 0) arc (0 : \deltaDeg : 0.22);
    \node at (\hingeX + 0.38, -0.09) {$\delta(t)$};
\end{scope}

% ──────────────────────────────────────────────────────────────────
% Quote d_x between EA and CG_f
\draw[quotestyle] (\EAx, -0.72) -- (0.87, -0.72);
\node[below] at (0.585, -0.72) {$d_{x}$};
\draw[dotted, thin, gray!70] (\EAx, 0) -- (\EAx, -0.72);
\draw[dotted, thin, gray!70] (0.87, -0.12) -- (0.87, -0.72);

% ──────────────────────────────────────────────────────────────────
% GUST W(t) — freccia verticale che scende dall'alto verso l'ala
% Posizionata sopra il profilo alare, al centro della corda
\draw[flowarrow] (-0.20, -0.50) -- (-0.20, -0.2);
\node[above] at (-0.20, -0.70) {$W(t)$};

% ─── Reference frame x-y ─────────────────────────────────────
\coordinate (frameorigin) at (-0.90, -0.70);
\draw[arrowstyle] (frameorigin) -- ++(0.22, 0) node[right] {$x$};
\draw[arrowstyle] (frameorigin) -- ++(0, -0.22) node[below] {$y$};
\fill (frameorigin) circle (0.012);

\end{tikzpicture}